% DELTA23: /root/pentagon, 2026-11-21 CET. Full proof draft.
\subsection{A small affine evaluation detecting the cubic frame sign}

\paragraph{D23.P3 (an order-1944 free-group detector).}
Put $k=\mathbb F_9=\mathbb F_3[i]$, $i^2=-1$, and use affine
multiplication $(v,M)(w,N)=(v+Mw,MN)$ on $k^2\rtimes\operatorname{GL}_2(k)$.
Retain the modular generators $s,r$, $x=(sr)^{-2}$, $y=sxs$ of
LM7\_PUSH E.58 and E.65. Set
\[
 S=\begin{pmatrix}0&1\\-1&0\end{pmatrix},\quad
 R=\begin{pmatrix}0&-1\\1&-1\end{pmatrix},\quad
 U=\begin{pmatrix}1&1\\0&1\end{pmatrix},\quad v=e_1.
\]
There is a continuous homomorphism
\[
 \eta:\widehat B=C_2\amalg C_3\longrightarrow E,
 \qquad \eta(s)=(0,iS),\qquad \eta(r)=(v,R),
 \tag{D23.P3.1}
\]
onto the finite group
\[
 E=k^2\rtimes H,\qquad H=\langle iI\rangle\operatorname{SL}_2(\mathbb F_3),
 \quad |E|=3888.
\]
Its restriction to the pure free group $F=\widehat F(x,y)$ has
image $E_F=k^2\rtimes\operatorname{SL}_2(\mathbb F_3)$, of order
$1944$. The image of $C_E(\eta(r))$ in $H$ is exactly $\langle R\rangle$,
and its order is $27$.

For an actual multiplier-one gentle element $h=(1,f)$, write
\[
 \eta(f)=(d,M_\rho),\quad d=(d_1,d_2)^{\mathsf t},\quad
 M_\rho=\begin{pmatrix}1&\rho\\\rho&1+\rho^2\end{pmatrix},
 \qquad\rho\in\mathbb F_3.
\]
Here $\rho$ is the reduction of the frame parameter in E.65; it is
also read directly as the lower-left entry of the linear part.
Then
\[
 \boxed{\quad (-1)^{a_3(h)}
       =(1-\rho)(1-d_1)-d_2\quad\text{in }k.\quad}
 \tag{D23.P3.2}
\]
The right-hand side belongs to $\{1,-1\}\subset\mathbb F_3$ for
every actual gentle element. Moreover, if $\epsilon=(-1)^{a_3(h)}$,
its full affine evaluation is exactly
\[
 d_1=\frac{1-\rho-\epsilon}{1+i},\qquad
 d_2=(\rho+i)d_1.
 \tag{D23.P3.3}
\]
Thus every such $f$-evaluation lies in an explicitly specified set of six elements of $E_F$.
For the E.58 witness one obtains
\[
 \eta(f_{58})=((1-i,1+i)^{\mathsf t},I),
 \tag{D23.P3.4}
\]
not merely the value of its sign.

\begin{proof}
The matrix $iS$ squares to $I$. In characteristic three,
$(R-I)^2=0$, so $I+R+R^2=0$. These identities prove the two
torsion relations in (D23.P3.1). Hence the homomorphism exists
by the defining free-product presentation, and continuity follows
because its target is finite.

The linear image is $H$: indeed $iSR=iU^{-1}$ and
$(iU^{-1})^3=-iI$, so it contains $iI,S,R$, which generate $H$.
The intersection $\langle iI\rangle\cap\operatorname{SL}_2(\mathbb F_3)$
is $\{I,-I\}$; hence $|H|=4\cdot24/2=48$.
The translation subgroup $N$ of the image is an $H$-invariant
$\mathbb F_3$-subspace of $k^2$. Scalar multiplication by $iI$
makes it a $k$-subspace. The natural representation of
$\operatorname{SL}_2(\mathbb F_3)$ over $k$ is irreducible:
a common invariant line for the two nonidentity unipotent
matrices $U$ and $L=\left(\begin{smallmatrix}1&0\\-1&1\end{smallmatrix}\right)$
would have to be both $ke_1$ and $ke_2$. Thus $N$ is zero or $k^2$.

If $N=0$, the affine image is the graph of a cocycle $c:H\to k^2$.
The central element $z=-I$ acts as $-1$. Comparing $c(gz)$ and
$c(zg)$ gives $2c(g)=(1-g)c(z)$, so
$c(g)=(1-g)t$ with $t=c(z)/2$. In particular this would imply
$v=c(R)\in(1-R)k^2$. But $(1-R)k^2=k(1,-1)^{\mathsf t}$
and $v=e_1$ is outside that line. This contradiction proves
$N=k^2$ and the asserted full image.

The linear images of $x,y$ are $-U^2,-L^2$. Cubing either
produces $-I$, so they generate $\operatorname{SL}_2(\mathbb F_3)$.
The pure image is normal in $E$, since $F$ is normal in $B$.
Its intersection with $k^2$ is again zero or all of $k^2$.
The zero alternative would make the quotient of $E$ by the pure
image have order $3888/24=162$, whereas it is a quotient of
$B/F\cong S_3$. It follows that the pure image is exactly $E_F$.

An element of the centralizer of $R$ in $H$ has the form
$i^jR^m$, $j\in\mathbb Z/4$, $m\in\mathbb Z/3$.
Indeed $C_{\operatorname{SL}_2(\mathbb F_3)}(R)
=\langle-I\rangle\times\langle R\rangle$, and the scalar
extension adds $iI$. For an affine element $(t,i^jR^m)$ to
centralize $(v,R)$ one must have
$(1-R)t=v-i^jR^m v$. In the one-dimensional coinvariant quotient
$k^2/(1-R)k^2$, the nonzero class of $v$ is fixed by $R$.
Consequently $i^j=1$. Conversely all powers of $(v,R)$ centralize
it, so the linear image of its centralizer is exactly
$\langle R\rangle$. Each of its three linear values has
$|\ker(1-R)|=9$ translation lifts. This proves the order $27$.

We next justify the signed comparison between the linear part of
$\eta$ and the supplied $q_{12}$, rather than identifying the
representations without a sign convention. There is a character
$\chi_4:\operatorname{SL}_2(\mathbb Z/4)\to\langle i\rangle$
with $\chi_4(U)=\chi_4(L)=-i$, $\chi_4(S)=i$,
$\chi_4(R)=1$, and $\chi_4(-I)=-1$. Its finite definition is
verified by enumerating all $48$ matrices and propagating the
values along multiplication by $U,L$; the certificate checks
consistency on every such edge. Thus
\[
 [(g_4,g_3)]\longmapsto\chi_4(g_4)g_3
\]
is a well-defined homomorphism from
$(\operatorname{SL}_2(\mathbb Z/4)\times
\operatorname{SL}_2(\mathbb F_3))/\langle(-I,-I)\rangle$ to $H$.
It sends the modular $s,r$ to $iS,R$, precisely the linear part
of $\eta$. For the actual pure frame $A_h$, E.65 proves
$q_4(A_h)=I$, so its linear image under $\eta$ is exactly
$q_3(A_h)$. For $f\in F'$ the modulo-four coordinate is also
trivial. E.65.9 therefore gives its displayed linear part
$M_\rho$, and $V=U^{2\rho}$ is the reduction of E.65's shear.

Write $\eta(A_h)=(t,A_0)$ and $T=A_0^{-1}V$. By E.65,
$T\in\langle-I\rangle\times\langle R\rangle$, and its sign
on the fixed line is $\epsilon=(-1)^{a_3(h)}$. The exact frame
identity $h(r)=A_hrA_h^{-1}=f^{-1}r$ gives, on translations,
\[
 M_\rho^{-1}(v-d)
       =A_0v+(1-A_0RA_0^{-1})t.
\]
Multiplying by $V^{-1}$ and passing to $R$-coinvariants makes
the last term zero. The row vector $\ell=(1,1)$ realizes this
coinvariant quotient; it satisfies $\ell R=\ell$, $\ell v=1$,
and $\ell T^{-1}=\epsilon\ell$. We obtain
\[
 \ell V^{-1}M_\rho^{-1}(v-d)=\epsilon.
\]
Direct multiplication gives
$V^{-1}M_\rho^{-1}=\left(\begin{smallmatrix}1&0\\-\rho&1\end{smallmatrix}\right)$,
which proves (D23.P3.2).

Finally $h(s)=sf$ has order two. The translation part of
$\eta(sf)^2=1$, with the displayed matrices, is equivalent to
$d_2=(\rho+i)d_1$. Substitution in (D23.P3.2) gives
$\epsilon=1-\rho-(1+i)d_1$, proving (D23.P3.3).
E.58 has $\rho=0$ and $a_3=1$, which gives (D23.P3.4).
The exact finite program \path{agents/pentagon/affine_cubic_detector.py}
independently enumerates both images, the centralizer, and all
six affine evaluations, and checks the coinvariant calculation
for all $486$ combinations of sign, cubic power, and translation.
\end{proof}

\paragraph{D23.P4 (an explicit finite evaluation gate for A.84).}
Retain $p=(123)$ and $t=(1234)$ in $S_4$. The simultaneous map
\[
 F\longrightarrow E_F\times S_4,\qquad
 x\longmapsto(\eta(x),p),\quad
 y\longmapsto(\eta(y),t)
 \tag{D23.P4.1}
\]
is onto a group of order $46656$. For an actual gentle element
$h=(1,f)$, write its $S_4$-evaluation as
$f(p,t)=t^{2\nu_4}p^{a_f}$, as in DELTA\_ON\_GT FP.37.
For actual full GT elements, the existing pentagon consequences
FP.33 and FP.37 give $a_f=\rho$ and $a_2=\nu_4$.
Consequently A.84 is equivalent to the following concrete finite
comparison on actual full-GT source residues:
\[
 \boxed{\quad
 (1-\rho)(1-d_1)-d_2=(-1)^{\nu_4}.
 \quad} \tag{D23.P4.2}
\]
After the two previously proved pentagon consequences, there are
only twelve possible joint evaluations: $\rho\in\mathbb F_3$
and $(a_3,\nu_4)\in\mathbb F_2^2$. Six obey (D23.P4.2) and
six fail it. The remaining question is whether any of the latter
six is the image of an actual full-GT element. This is not an
assertion that any particular finite candidate lifts to GT.

\begin{proof}
Since $-I$ acts as $-1$ on $k^2$ and $2$ is invertible, commutators
with $-I$ generate the translation subgroup. Hence
$E_F^{\mathrm{ab}}\cong\operatorname{SL}_2(\mathbb F_3)^{\mathrm{ab}}
\cong C_3$. On the other hand $S_4^{\mathrm{ab}}\cong C_2$.
Both groups are finite solvable. Any nontrivial common quotient
would be solvable and would have a nontrivial abelian quotient;
its abelianization would be a quotient of both $C_3$ and $C_2$,
which is impossible. Both projections of (D23.P4.1) are onto,
so Goursat's lemma proves that its image is the full product.
Its order is $1944\cdot24=46656$.

D23.P3 computes the cubic bit from its first coordinate, while
FP.37 computes the harmonic bit on full GT from its second
coordinate. Thus their equality is exactly (D23.P4.2).
The six possible first coordinates are (D23.P3.3); the second
coordinates have the six values $t^{2\nu_4}p^{a_f}$.
Imposing $a_f=\rho$ leaves the stated twelve candidates and
partitions them according to the displayed sign equality.
The reduction from arbitrary multiplier to multiplier one is
LM7\_PUSH E.65.6, using Galois vanishing, exactly as in
DELTA\_ON\_GT AO.1. This proves the equivalence with A.84.
\end{proof}
